\documentclass[12pt,a4paper]{article}

\usepackage[utf8]{inputenc}
\usepackage[spanish]{babel}
\usepackage{amsmath}
\usepackage{amsfonts}
\usepackage{amssymb}
\usepackage{geometry}
\geometry{a4paper, margin=1in}

\begin{document}

\begin{flushleft}
\textbf{Arquitectura del Computador} \hfill \textbf{Freddy Ovalle} \\
\textbf{Prof. José Canache} \hfill \textbf{C.I: 31.709.650} \\
\textbf{Sección 01}
\end{flushleft}

\vspace{1cm}

\begin{center}
\Large{\textbf{Fórmulas de Cálculo de Prestaciones}}
\end{center}

\vspace{0.5cm}

\section*{Prestaciones relativas}

\begin{itemize}
    \item \textbf{Definición de Prestaciones:} 
    $$Prestaciones = \frac{1}{\text{Tiempo de Ejecución}}$$

    \item \textbf{Rendimiento Relativo:} 
    $$n = \frac{Prestaciones_X}{Prestaciones_Y} = \frac{Tiempo_Y}{Tiempo_X}$$
\end{itemize}

\section*{Ecuación Clásica de la CPU}

\begin{itemize}
    \item \textbf{Tiempo de CPU (Ciclos):} 
    $$\text{Tiempo de CPU} = \text{Ciclos de Reloj de la CPU} \times \text{Tiempo de Ciclo de Reloj}$$

    \item \textbf{Tiempo de CPU (Frecuencia):} 
    $$\text{Tiempo de CPU} = \frac{\text{Ciclos de Reloj de la CPU}}{\text{Frecuencia de Reloj}}$$

    \item \textbf{Ciclos de Reloj Totales:} 
    $$\text{Ciclos de Reloj de la CPU} = \text{Número de Instrucciones} \times CPI$$

    \item \textbf{Ecuación Completa (Tiempo):} 
    $$\text{Tiempo de CPU} = \text{Número de Instrucciones} \times CPI \times \text{Tiempo de Ciclo}$$

    \item \textbf{Ecuación Completa (Frecuencia):} 
    $$\text{Tiempo de CPU} = \frac{\text{Número de Instrucciones} \times CPI}{\text{Frecuencia de Reloj}}$$

    \item \textbf{Ciclos de CPU por Mezcla de Instrucciones:} 
    $$\text{Ciclos de CPU} = \sum (CPI_i \times C_i)$$

    \item \textbf{CPI Medio:} 
    $$CPI = \frac{\text{Ciclos de Reloj de la CPU}}{\text{Número de Instrucciones}}$$
\end{itemize}

\section*{Segmentación (Pipelining)}

\begin{itemize}
    \item \textbf{Tiempo Ideal entre Instrucciones:} 
    $$\text{Tiempo entre Instrucciones (Seg)} = \frac{\text{Tiempo entre Instrucciones (No Seg)}}{\text{Número de Etapas}}$$
\end{itemize}

\section*{Ley de Amdahl}

\begin{itemize}
    \item \textbf{Tiempo de Ejecución Mejorado:} 
    $$\text{Tiempo Post-Mejora} = \left(\frac{\text{Tiempo Afectado}}{\text{Factor de Mejora}}\right) + \text{Tiempo No Afectado}$$
\end{itemize}

\section*{Memoria y Caches}

\begin{itemize}
    \item \textbf{Ciclos Parada de Memoria (Accesos):} 
    $$\text{Ciclos Parada} = \left(\frac{\text{Accesos a Memoria}}{\text{Programa}}\right) \times \text{Frecuencia Fallos} \times \text{Penalización Fallo}$$

    \item \textbf{Ciclos Parada de Memoria (Instrucciones):} 
    $$\text{Ciclos Parada} = \left(\frac{\text{Instrucciones}}{\text{Programa}}\right) \times \left(\frac{\text{Fallos}}{\text{Instrucción}}\right) \times \text{Penalización Fallo}$$

    \item \textbf{Ciclos Parada por Lectura:} 
    $$\text{Ciclos Parada Lectura} = \left(\frac{\text{Lecturas}}{\text{Programa}}\right) \times \text{Frecuencia Fallos Lectura} \times \text{Penalización Lectura}$$

    \item \textbf{Ciclos Parada por Escritura:} 
    $$\text{Ciclos Parada Escritura} = \left(\left(\frac{\text{Escrituras}}{\text{Programa}}\right) \times \text{Frecuencia Fallos Escritura} \times \text{Penalización}\right) + \text{Parada Búfer}$$
    
    \item \textbf{Tiempo Medio de Acceso a Memoria (AMAT):}
    $$AMAT = \text{Tiempo Acierto} + (\text{Frecuencia Fallos} \times \text{Penalización Fallo})$$
\end{itemize}

\section*{Métricas Alternativas}

\begin{itemize}
    \item \textbf{MIPS (Tiempo de Ejecución):}
    $$MIPS = \frac{\text{Número de Instrucciones}}{\text{Tiempo de Ejecución} \times 10^6}$$
    
    \item \textbf{MIPS (Reloj y CPI):}
    $$MIPS = \frac{\text{Frecuencia de Reloj}}{CPI \times 10^6}$$
\end{itemize}

\end{document}
